\clearpage
\VerdierPhysicalPageAnchor{0130}%
\begingroup
\small
\setlength{\emergencystretch}{1.5em}%
\setlength{\abovedisplayskip}{2pt plus 1pt minus 1pt}%
\setlength{\belowdisplayskip}{2pt plus 1pt minus 1pt}%

\VerdierSemanticAnchor{ch2-2-1-3-2-sm2-completion}
\noindent can be completed to a commutative diagram:

\VerdierSemanticAnchor{ch2-2-1-3-2-completion-diagram}
\begin{center}
  $t\in S$\qquad
  \begin{tikzcd}[column sep=4.0em,row sep=2.5em,baseline=(current bounding box.center)]
    Z' \arrow[r,dashed,"g"] \arrow[d,dashed,"t"'] &
    X' \arrow[d,"s"] \\
    Z \arrow[r,"f"'] & X
  \end{tikzcd}
  \quad ,
\end{center}

\VerdierSemanticAnchor{ch2-2-1-3-2-opposite-reduction}
\noindent for then the second assertion of SM2) will follow by passing to the
opposite category.

\VerdierSemanticAnchor{ch2-2-1-3-2-triangle-construction}
By TRI) and TRII), we know that there exists a distinguished triangle
$(X,Y,Z[1],u,v,f[1])$ and, by TRI), we know that there exists a distinguished
triangle $(X',Y,Z',us,v',w')$. We then have a commutative diagram:

\VerdierSemanticAnchor{ch2-2-1-3-2-us-square}
\begin{center}
  \begin{tikzcd}[column sep=4.2em,row sep=2.5em]
    X' \arrow[r,"us"] \arrow[d,"s"'] &
    Y \arrow[d,"\operatorname{id}_Y"] \\
    X \arrow[r,"u"'] & Y
  \end{tikzcd}
  \quad .
\end{center}

\VerdierSemanticAnchor{ch2-2-1-3-2-sm6-application}
\noindent The morphism $\operatorname{id}_Y$ belongs to $S$ (SM1) and, by
SM6), we obtain a morphism $s'\colon Z'\to Z[1]$ belonging to $S$ such that
$(s,\operatorname{id}_Y,s')$ is a morphism of triangles. In particular, the
following diagram is commutative:

\VerdierSemanticAnchor{ch2-2-1-3-2-shift-square}
\begin{center}
  \begin{tikzcd}[column sep=4.2em,row sep=2.5em]
    Z' \arrow[r,"{w'}"] \arrow[d,"{s'}"'] &
    X'[1] \arrow[d,"{s[1]}"] \\
    Z[1] \arrow[r,"{f[1]}"'] & X[1]
  \end{tikzcd}
  \quad ,
\end{center}

\VerdierSemanticAnchor{ch2-2-1-3-2-proof-complete}
\noindent whence, using SM5), the first assertion of SM2).

\VerdierSemanticAnchor{ch2-2-1-4}
\medskip
\noindent\textbf{2.1.4.}\quad
\VerdierSemanticAnchor{ch2-2-1-4-convention}
\emph{Unless expressly stated otherwise, the multiplicative systems of
triangulated categories will be compatible with the triangulation. They will
therefore be called simply, and by abuse of language, multiplicative systems.}

\VerdierSemanticAnchor{ch2-2-1-5}
\medskip
\noindent\textbf{2.1.5.}\quad
\VerdierSemanticAnchor{ch2-2-1-5-full-triangulated-subcategory}
Let $\mathcal D$ be a triangulated category. A \emph{full triangulated
subcategory} of $\mathcal D$ is a full subcategory of $\mathcal D$ stable
under translations, equipped with a triangulated structure such that the
inclusion functor is an exact functor (for strict commutation with
translations).

\VerdierSemanticAnchor{ch2-2-1-5-additive-subcategory}
A full triangulated subcategory of $\mathcal D$ is an additive subcategory of
$\mathcal D$ (1.2.7).

\VerdierSemanticAnchor{ch2-2-1-5-boundary}
A full subcategory
\endgroup
